%---------------------------------------------------------------------
\chapter{Agent Design and Architecture}
\label{ch:design}
%---------------------------------------------------------------------

\section{System Architecture}
Figure~\ref{fig:arch} shows the architecture. The design separates four concerns that the failure modes of
a compliance assistant make dangerous to mix: storing evidence, computing facts, deciding what to
investigate, and deciding what to accept. Evidence is stored once, immutably and by digest. Facts are
computed by deterministic code. The model decides only which check to run and how to read the result.
A gate and a verifier decide what the system is willing to say, and a human decides the report.

\begin{figure}[H]
\centering
\resizebox{\linewidth}{!}{%
\begin{tikzpicture}[
  box/.style={draw=bmsblue,fill=guidebg,rounded corners=2pt,minimum width=2.9cm,minimum height=1.0cm,
              align=center,font=\footnotesize},
  det/.style={draw=ForestGreen,fill=dobg,rounded corners=2pt,minimum width=2.9cm,minimum height=1.0cm,
              align=center,font=\footnotesize},
  llm/.style={draw=Orange,fill=Orange!12,rounded corners=2pt,minimum width=2.9cm,minimum height=1.0cm,
              align=center,font=\footnotesize},
  store/.style={draw=gray,fill=gray!10,cylinder,shape border rotate=90,aspect=0.2,minimum width=2.9cm,
              minimum height=1.2cm,align=center,font=\footnotesize},
  human/.style={draw=BrickRed,fill=dontbg,rounded corners=2pt,minimum width=2.9cm,minimum height=1.0cm,
              align=center,font=\footnotesize},
  arr/.style={-{Stealth[length=2mm]},thick},
  lab/.style={font=\scriptsize,fill=white,inner sep=1pt}]
% Row 1: evidence in
\node[box] (files) at (0,0) {Dossier files\\PDF, CSV, JSON};
\node[det] (intake) at (4.9,0) {Intake and\\extraction};
\node[det] (norm) at (9.8,0) {Normalisation,\\frozen snapshots};
\node[store] (blob) at (14.7,0) {Blob store,\\audit kernel};
% Row 2: knowledge and facts
\node[box] (corpus) at (0,-2.6) {Regulatory corpus\\Rules 2022, FAQ};
\node[det] (index) at (4.9,-2.6) {FTS5 index,\\retrieval grader};
\node[det] (tools) at (9.8,-2.6) {Audited read-only\\tools (6)};
\node[det] (recon) at (14.7,-2.6) {Exact reconciler\\(decimals)};
% Row 3: agent
\node[det] (gate) at (0,-5.2) {Finding gate\\citations, digits, words};
\node[llm] (model) at (4.9,-5.2) {Model gateway\\Qwen3.5-9B 4-bit};
\node[det] (graph) at (9.8,-5.2) {LangGraph runner\\preflight, investigate,\\draft, verify};
% Row 4: output
\node[det] (report) at (4.9,-7.8) {Report builder\\and verifier};
\node[human] (review) at (9.8,-7.8) {Human reviewer\\decision};
\node[det] (export) at (14.7,-7.8) {Export, replay,\\evaluation};
\draw[arr] (files) -- (intake);
\draw[arr] (intake) -- (norm);
\draw[arr] (norm) -- (blob);
\draw[arr] (corpus) -- (index);
\draw[arr] (index) -- node[lab,above]{provisions} (tools);
\draw[arr] (recon) -- node[lab,above]{balances} (tools);
\draw[arr] (norm.south east) -- node[lab,sloped,above]{snapshot} (recon.north west);
\draw[arr] (graph) -- node[lab,right]{tool call} (tools);
\draw[arr,<->] (graph) -- node[lab,above]{PLAN, ASSESS} (model);
\draw[arr] (model) -- node[lab,above]{proposal} (gate);
\draw[arr] (gate.south east) -- node[lab,sloped,above]{accepted} (report.north west);
\draw[arr] (graph.south west) -- node[lab,sloped,above]{draft} (report.north east);
\draw[arr] (report) -- node[lab,above]{to review} (review);
\draw[arr] (review) -- node[lab,above]{bound decision} (export);
\draw[arr,dashed] (blob.east) -- ++(1.0,0) |- node[lab,pos=0.25,right]{append-only events} (export.east);
\end{tikzpicture}}
\caption{Architecture of ReCycleGuard. Green boxes are deterministic, the orange box is the language
model, the red box is the human decision, and the cylinder is durable evidence.}
\label{fig:arch}
\end{figure}

\section{Component Description}

\subsection{Intake, Extraction and Normalisation}
Intake receives untrusted files, hashes their bytes, writes them to a content-addressed blob store and
issues a receipt that records the outcome, including failures. Extractors turn PDF pages, CSV rows and JSON
records into \emph{observations}, each with a precise locator. A PDF page with no extractable text is
recorded as incomplete coverage and is not silently skipped. Normalisation maps units and entity aliases to
canonical forms. A failure at any of these stages is retained as an outcome and never dropped, because
a missing record is evidence in itself.

\subsection{Snapshots and Calculation Profile}
A case snapshot and a corpus snapshot are frozen with digests, and the calculation profile for the period is
pinned. A run records the exact digests it used. An investigation therefore cannot see a later edit of the
case, and a reviewer can tell which version of the evidence a report was built from.

\subsection{Retrieval and Relevance Grader}
The corpus is indexed with SQLite FTS5 and ranked with BM25. Retrieval proceeds in stages: eligibility
filtering by regime, role and reviewed temporal assertions; candidate search; bounded expansion to parent
headings, footnotes and amendment counterparts; grading; and at most two query reformulations per finding.
The grader records authority, temporal applicability, direct support, completeness and contradiction
independently. Passages whose applicability cannot be established are labelled as leads and cannot support a
definitive statement. If retrieval is weak or contradictory, the system escalates to a human and does not
guess.

\subsection{Exact Reconciler}
For each obligation bucket, defined by producer, liability year, metal and unit, the reconciler computes
\begin{equation}
\text{gap}_b = \text{required}_b - \text{supported}_b,
\end{equation}
where $\text{supported}_b$ is the sum of active allocations into bucket $b$ derived from the allocation
ledger. Pending quantities are shown separately and are not subtracted. Metal gaps are never added into a
legally meaningful total. A supplied portal balance is stored as a comparison-only observation, and the
reconciler reports the delta between the portal figure and the ledger-supported figure instead of
treating either as the truth. All values are decimal strings and the arithmetic is exact.

\subsection{Tools}
The model reaches the case only through six tools, listed in Table~\ref{tab:tools} in Chapter~\ref{ch:impl}.
All are read-only. A request for any other operation aborts the run. Each call writes an intent record before
it executes and a result record after it completes.

\subsection{LangGraph Runner and Model Gateway}
The runner is a LangGraph state graph with the nodes \emph{preflight}, \emph{investigate}, \emph{draft},
\emph{verify} and \emph{sync}. Each call to the runner executes exactly one turn and commits it, so a process
that is stopped mid-run resumes from the last committed turn without repeating a tool call. The model gateway
is a client of the local MLX-LM server. It builds the prompt, sends one request, validates the reply against a
JSON schema, and records the exact request and response, the token counts and the latency.

\subsection{Finding Gate, Report Builder and Verifier}
A model may attach one finding proposal to an assessment. The gate accepts it only if every citation is
evidence that this run's tools returned, the text contains no digits and no prohibited or authority wording,
and a ``clean'' finding is supported by a reconciliation from this run with nothing pending. Rejections are kept
in the report with their reasons. The report builder assembles the structured report before any prose is
rendered, and the verifier checks that citations exist and resolve, that the numbers agree with the
calculations, and that the scope language is within bounds. After one failed repair the report is blocked.

\subsection{Human Review, Export and Replay}
The review service presents a packet to an authenticated reviewer and binds the decision to the exact report
and input digests. Export produces a bundle of the complete audit trail, which a replay port can use to
serve recorded model replies by request digest. A replay is not a new trial and is not evidence of model
quality.

\section{Selected Agentic Patterns and Justification}
\label{sec:patterns}
\begin{table}[H]
\centering
\caption{Patterns used, alternatives considered and the reason for the choice}
\label{tab:patterns}
\small
\begin{tabularx}{\linewidth}{@{}p{3.1cm} X X@{}}
\toprule
\textbf{Pattern} & \textbf{How it is used} & \textbf{Reason and alternative rejected} \\ \midrule
Plan and assess loop & The model proposes one check, a tool runs it, the model assesses the observation, and the loop continues, resolves or abstains. & Gives visible adaptation and a record of each step. Free-form ReAct was rejected because its steps are hard to bound and validate. \\
Tool use & Six whitelisted read-only tools with typed arguments. & Keeps arithmetic and lookup out of the model. A code-execution tool was rejected as unbounded. \\
Corrective retrieval & Graded BM25 retrieval with at most two reformulations and escalation. & Authority and date matter as much as relevance in legal text. Dense retrieval is optional in the design and is not part of the evaluated system. \\
Model proposes, code disposes & The finding gate and the report verifier check every model-originated claim. & A small quantised model cannot be trusted to cite and calculate. Trusting model text was rejected. \\
Bounded loops and budgets & Limits on calls, tokens, tool calls per finding, reformulations, retries and time. & Prevents runaway cost and cycles; identical repeated actions are stopped by a fingerprint check. \\
Human in the loop & No report is final until a reviewer decides it. & The system is advisory by design. Autonomous approval was excluded. \\
Deterministic preflight & Coverage and profile checks run before any model call. & Missing pages must stop the run before the model can guess around them. \\
\bottomrule
\end{tabularx}
\end{table}

Two investigation variants share the same tools, corpus and budgets, and they differ in who chooses the
next check. In the fixed-retrieval variant (V1) a pre-declared sequence of checks runs and the model only
assesses. In the adaptive variant (V2) the model chooses each check from the allowed tools. The comparison of
the two is the experiment of Chapter~\ref{ch:exp}.

\section{Control Flow and Agent State}
\label{sec:flow}
Figure~\ref{fig:states} shows the run states. The run starts with a deterministic preflight. If coverage is
incomplete, for example because a page has no extractable text, the run stops in a blocked-input state before
any model call. Otherwise the investigation alternates between a planning call and an assessment call.
Each assessment ends with a disposition: continue, resolve or abstain. A resolve disposition requires the
verdict to support the finding, and any other verdict forces continue or abstain. When no finding remains
pending, or the plan stops with no new evidence, the run drafts a report. A failed verification returns
once to drafting and then blocks the output.

\begin{figure}[H]
\centering
\resizebox{\linewidth}{!}{%
\begin{tikzpicture}[
  st/.style={draw=bmsblue,fill=guidebg,rounded corners=3pt,minimum width=3.0cm,minimum height=0.9cm,
             align=center,font=\footnotesize},
  term/.style={draw=gray,fill=gray!10,rounded corners=3pt,minimum width=3.0cm,minimum height=0.9cm,
             align=center,font=\footnotesize},
  bad/.style={draw=BrickRed,fill=dontbg,rounded corners=3pt,minimum width=3.0cm,minimum height=0.9cm,
             align=center,font=\footnotesize},
  arr/.style={-{Stealth[length=2mm]},thick},
  lab/.style={font=\scriptsize,fill=white,inner sep=1pt}]
\node[st] (pre) at (0,0) {PREFLIGHT};
\node[st] (inv) at (5.4,0) {INVESTIGATING\\(plan, act, assess)};
\node[st] (dra) at (10.8,0) {DRAFTING};
\node[st] (ver) at (16.2,0) {VERIFYING};
\node[bad] (bin) at (0,-2.3) {BLOCKED\_INPUT};
\node[bad] (bout) at (10.8,-2.3) {BLOCKED\_OUTPUT};
\node[bad] (fail) at (5.4,-2.3) {FAILED};
\node[st] (await) at (16.2,-2.3) {AWAITING\_REVIEW};
\node[term] (done) at (16.2,-4.3) {COMPLETED\\(reviewer decides)};
\draw[arr] (pre) -- node[lab,above]{coverage ok} (inv);
\draw[arr] (inv) -- node[lab,above]{none pending} (dra);
\draw[arr] (dra) -- (ver);
\draw[arr] (pre) -- node[lab,right]{page missing} (bin);
\draw[arr] (inv) -- node[lab,right]{budget, provider error} (fail);
\draw[arr] (dra) -- node[lab,right]{blocked} (bout);
\draw[arr] (ver) -- node[lab,right]{verified} (await);
\draw[arr] (ver.north) -- ++(0,0.7) -| node[lab,pos=0.25,above]{repair once} (dra.north);
\draw[arr] (ver.south west) -- node[lab,pos=0.55,above,sloped]{repairs exhausted} (bout.north east);
\draw[arr] (await) -- (done);
\draw[arr] (inv.north) -- ++(0,0.7) -| node[lab,pos=0.25,above]{continue} (inv.north east);
\end{tikzpicture}}
\caption{Run states of the investigation workflow}
\label{fig:states}
\end{figure}

\textbf{Agent state.} The persisted state consists of the run identifier, a version number, the pinned
snapshots and profile, the list of findings with their status, the budget usage, and the sequence of turn
manifests. Every update uses compare-and-swap on the version number, so two writers cannot overwrite each
other. The fingerprint of each proposed action, made of the finding, the tool and the arguments, is checked
against earlier actions, and an exact repeat stops the run with a cycle outcome.

\textbf{Budgets.} All real-model runs use one frozen budget: at most 12 model calls, 600\,s of wall time
and a token reservation per call that is checked before dispatch. The same limits apply to both variants so
that the comparison is fair.

\section{Data Flow}
\label{sec:dataflow}
Evidence flows in one direction. Original bytes enter the blob store and are never changed. Extractors derive
observations from them, normalisation derives canonical records, and the frozen snapshots collect those into
a case version. The reconciler reads the snapshot and emits calculations. A tool call reads the snapshot or the
index and returns a typed observation with identifiers. The model sees these identifiers and may cite only
identifiers that appear in them. The gate and the verifier check each citation against the tool outputs of the
same run. Finally the report, the review decision and the full event stream are stored as linked records.
Nothing flows back from the model into the evidence: a model output can change a report only by passing
through a typed proposal and the gate.
